\section{A Generic DKG Construction} \label{generic-construction}

We now introduce a weakly robust generic construction that assumes an honest majority, shown in Figure~\ref{fg:generic-perfect}.
We present a concrete constructions in Section~\ref{storm-gen}.

\begin{figure}[p]
    % chelsea- make full size for eprint
\begin{pchstack}[boxed,center,space=1em]
\begin{pcvstack}
	%\pseudocode[codesize=\scriptsize]{
\procedure[linenumbering]{$\dkgkeygen_0[\Hash](\lambda, n, t, i)$} {
    (\secretone_i, A_i) \rgets \enike.\keygen(\lambda) \\
    \mathsf{in} \gets (\lambda, \secretone_i, n, t) \\
    \{ (j, \share_{ij}) \}_{j =1}^n, \vsscommit_i \rgets \avss[\Hash_1].\issueshares(\mathsf{in}) \\
    \text{\algcom The \ENIKE secret key $\secretone_i$ is} \\
    \text{\algcom secret-shared via the \AVSS.} \\
    (\witness_i, \commitment_i) \rgets \enike.\keygen(\lambda) \\
    %\text{\algcom Broadcast commitments} \\
    \bmsg_{0,i} \gets (A_i, \commitment_i, \vsscommit_i) \\
    \pcfor j \in \set{n}, j \neq i \pcdo \\
    %\text{\algcom Send shares via} \\
    %\text{\algcom peer-to-peer channel} \\
    \pcind \pmsg_{0,ij} \gets (j, \share_{ij}) \\
    \party_i \gets ( \witness_i, \share_{ii}, \vsscommit_i, \commitment_i  ) \\
  \outvecp \gets \{ \pmsg_{0,ij} \}_{j \in \set{n},j\neq i} \\
  \pcreturn (\party_i, \mathsf{out}_p,  \bmsg_{0,i})
}
\pcvspace
\procedure[linenumbering]{$\dkgkeygen_1[\Hash](\party_i, \invecp_1, \invecb_1)$} {
  \parse \{(i, \share_{ji}) \}_{j =1}^n \gets \invecp_1 \\
  \parse \{ (A_j, \commitment_j, \vsscommit_j) \}_{j =1}^n \gets \invecb_1 \\
  %
  \pcfor j \in \set{n}, j \neq i \pcdo \\
  \pcind \pcif \avss[\Hash_1].\verify(i, \share_{ji}, \vsscommit_j) \neq 1 \\
  \pcind \textbf{or } \avss[\Hash_1].\derivepshare(0, \vsscommit_j) \neq A_j \\
  \pcind \pcind \party_i.\status = \abort \\
  \pcind \pcind \bmsg_{1,i} \gets \fail \\
  %\text{\algcom If any verification check fails,} \\
  %\text{\algcom abort} \\
  \pcind \pcind \pcreturn (\party_i, \bot, \bmsg_{1,i}) \\
  \party_i \gets \party_i \cup \{ (j, \share_{ji}) \}_{j \in \set{n}} \\
  \bmsg_{1,i} \gets \accept \\
  \pcreturn (\party_i, \bot, \bmsg_{1,i})
}
\pcvspace
\procedure[linenumbering]{$\dkgkeygen_2[\Hash](\party_i, \emptyset, \invecb_2)$} {
  \pcif \fail \in \invecb_2 \\
    \pcind \party_i.\status = \abort \\
    %\text{\algcom If any participant fails, abort} \\
    \pcind \pcreturn (\party_i,\ \bot,\ \bot) \\
  %
  \bmsg_{2,i} \gets \witness_i \\
  \pcreturn (\party_i,\ \bot,\ \bmsg_{2,i})
}
\end{pcvstack}
\pchspace
\begin{pcvstack}
\procedure[linenumbering]{$\finalize[\Hash](\party_i, \emptyset, \invecb_3)$} {
  \pcif \party_i.\status = \abort \\
  \pcind \pcreturn (\bot, \bot,\bot), (\abort, \bot) \\
  \parse \{ \witness_j \}_{j \in \set{n}, j \neq i} \gets \invecb_3 \\
  \qual \gets \emptyset;\ \corrupt \gets \emptyset \\
  \pcfor j \in \set{n}, j \neq i \pcdo  \\
    \pcind \pcif \enike.\verifykeypair(\witness_j, \commitment_j) = 1  \\
    \pcind \pcind \blinding_j \gets \enike.\getsharedkey(\witness_j, A_j) \\
  \text{\algcom Qualified participants send} \\
  \text{\algcom the correct opening $\witness_j$} \\
    \pcind \pcind \qual \gets \qual \cup \{ j \} \\
    \pcind \pcelse \corrupt \gets \corrupt \cup \{ j \} \\
    \pcif \corrupt \neq \emptyset \\
    \pcind \textbf{Broadcast } \{ \share_{ij} \}_{j \in \corrupt} \\
    \text{\algcom Broadcast shares for corrupt} \\
    \text{\algcom parties to all other participants} \\
    \pcind \textbf{Receive } \{ \share_{kj} \}_{j \in \corrupt, k \in \qual} \\
    \pcfor j \in \corrupt \pcdo \\
    \pcind \secretone_j \gets \avss[\Hash_1].\recover(t, \{(k,\share_{jk}) \}_{k \in \qual} ) \\
    \pcind \blinding_j \gets \enike.\getsharedkey(\secretone_j, \commitment_j)  \\
    %
    %\privaggregateset =  \{\share_{ji}\}_{j \in \set{n}} \\
    %\pubaggregateset = \{\vsscommit_j\}_{j \in \set{n}} \\
    %
    \tweak \gets \avss[\Hash_1].\gettweak(\{\vsscommit_j\}_{j =1}^n,
    \{\blinding_j \}_{j =1}^n ) \\
    \sk_i \gets \avss[\Hash_1].\aggregatepriv(i,  \{\share_{ji}\}_{j \in \set{n}}, \tweak) \\
    \avsscommit \gets \avss[\Hash_1].\aggregatepub(\{\vsscommit_j\}_{j \in \set{n}}, \tweak) \\
    \pk \gets \avss[\Hash_1].\derivepshare(0, \avsscommit) \\
    \text{\algcom Derive public keys for all} \\
    \text{\algcom remaining participants} \\
    \pcfor i \in \qual \pcdo \\
    \pcind \pk_i = \avss[\Hash_1].\derivepshare(i, \avsscommit)  \\
    \auxinfo \gets \{ \pk_i \}_{i \in \qual} \\
    \pcreturn (\PK, \qual, \auxinfo), (\accept, \SK_i)
    %
}
\end{pcvstack}
\end{pchstack}
    \caption{A generic DKG key generation protocol as defined in
    Section~\ref{dkg-reqs},
    initialized with a NIKE $\enike$ as defined in Section~\ref{enike},
    an \AVSS $\avss$ as defined in Section~\ref{asym-vss},
    and a hash function $\Hash$.
    Each $\dkgkeygen_r$ operation is defined with respect to participant $i$.
    The peer-to-peer channel inputs for round $r$ are $\invecp_r$ and outputs are $\outvecp_r$,
    composed of the peer-to-peer messages $\pmsg_j$.
    The broadcast channel inputs are $\invecb_r$
    composed of the broadcast messages $\bmsg_j$,
    sent by participant $j, j \in \set{n}$.
    }
  \label{fg:generic-perfect}
\end{figure}

Our generic construction employs the following building blocks:

\begin{itemize}
  \item A secure hash function $\Hash$, where $\Hash_1$ and $\Hash_2$ are domain-separated instances of $\Hash$, such that:
\begin{itemize}
  \item $\Hash_1: (\{\vsscommit_j \}_{j \in \set{n}}, \auxstring) \mapsto
    \tweak$: Accepts as input a tuple, where the first element in the tuple is a set of $n$ \AVSS commitments ordered canonically,
    and the second element is an auxiliary element $\auxstring$ as output by $\Hash_2$ below.
    The output is the tweak employed by the \AVSS aggregation algorithm.
  \item $\Hash_2: \{ \blinding_j \}_{j \in \set{n}} \mapsto \auxstring$: Accepts as input a set of $n$ elements,
  and outputs an auxiliary element $\auxstring$.
\end{itemize}
  \item An aggregatable verifiable secret sharing scheme
    $\avss[\Hash_1]~=~(\issueshares,\verify,\allowbreak\recover,\allowbreak\derivepshare, \aggregatepriv, \aggregatepub)$ that is zero-indexed.
  \item A non-interactive key exchange (NIKE) $\enike~=~(\keygen,\verifykeypair,\getsharedkey)$.
\end{itemize}

The secret and public domains of the \AVSS and NIKE must be the same as that of the target key generation scheme,
and the distribution over the secret domain for the NIKE must be the same as that of the target key generation scheme.
Further, we assume that the NIKE outputs secret and public keys with strictly a one-to-one relation.

In addition to these building blocks,
our construction requires that these primitives are algebraically composable.
Intuitively,
the secret key output from $\enike.\keygen$ should be a valid secret that can be shared via $\avss.\issueshares$,
and the \AVSS commitment should be verifiable with respect to the public \ENIKE key.
Towards this latter point,
we require Assumption~\ref{ass:avss-enike} to hold.

\begin{assumption} \label{ass:avss-enike}
  For any $(\sk, \pk) \rgets \enike.\keygen(\lambda)$,
  and any $(n, t) \in \mathbb{N}$ such that $n \geq t$;
  and any $C \subseteq \set{n}, \lvert C \rvert \geq t$,
  Equation~\ref{eq:avss-enike-rel} holds:
  \begin{align} \label{eq:avss-enike-rel}
    \begin{split}
      \enike.\verify(\sk, \pk) = 1, \text{ where} \\
      (\{ (i, \share_i) \}_{i \in \set{n}}), \vsscommit) \rgets \avss[\Hash].\issueshares(\lambda, \sk, n, t) \text{ and } \\
      \pk \gets \avss[\Hash_1].\derivepshare(0, \vsscommit) \text{ and } \\
      \sk \gets \avss[\Hash_1].\recover(t, \{(i, \share_i) \}_{i \in C})
    \end{split}
  \end{align}
\end{assumption}


\subsubsection{Key Generation Protocol.}
As shown in Figure~\ref{fg:generic-perfect}, the key generation protocol is
completed in three network rounds in the optimistic setting and four network
rounds in the worst-case setting.

First, in $\dkgkeygen_0$, each participant $i$ performs $\enike.\keygen$ twice;
first to generate a secret key $\secretone_i$ and its commitment $A_i$,
and second to generate a random value $\witness_i$ and its commitment $\commitment_i$.
%
Then, each participant $i$ performs an \AVSS secret sharing of $\secretone_i$,
outputting $n$ secret shares $\{ (j, \share_{ij}) \}_{j \in \set{n}}$,
and an \AVSS commitment $\vsscommit_i$.
Finally, each participant $i$ broadcasts its \ENIKE and \AVSS commitments to all other players via a broadcast message $\bmsg_{0,i}$.
Additionally, each participant $i$ sends one secret share to every other player $j \neq i$ via a peer-to-peer message $\pmsg_{0,ij}$,
keeping one share for itself.

In $\dkgkeygen_1$,
each participant $i$ receives shares $ \{(i, \share_{ji}) \}_{j \in \set{n}, j \neq i} $
from all other players via the peer-to-peer channel input $\invecp_1$
and commitments $\{ (A_j, \commitment_j, \vsscommit_j) \}_{j \in \set{n}, j \neq i} $
from all other players via the broadcast channel input $\invecb_1$.
Each participant then verifies the correctness of their received share with respect to its corresponding commitment.
Additionally, each participant $i$ verifies that the secret committed to in $\vsscommit_i$ is the same secret committed to in $A_i$,
by checking that $\avss.\derivepshare(0, \vsscommit_i) = A_i$.
If any check fails, the participant will broadcast a $\fail$ message,
and set their status to $\abort$.
Any participant that sets their status to $\abort$ results in that participant terminating (immediately exiting) the protocol.
Else, it will broadcast an $\accept$ message.

In $\dkgkeygen_2$,
participants receive as input status messages from all other participants via the broadcast channel input $\invecb_2$.
If any participant received a $\fail$ message,
that participant immediately exits $\dkgkeygen_2$
after setting their status to $\abort$\footnote{Because the fail message is sent on a broadcast channel,
if any participant receives a fail message, then all participants receive that message}.
Otherwise, each participant $i$ broadcasts $\witness_i$ to all other participants (which is the opening
to the commitment $\commitment_i$ sent in $\dkgkeygen_0$).

Finally, in $\finalize$, each participant $i$ receives as input the openings $\{ \witness_j \}_{j \in \set{n}, j \neq i} $
from all other participants via the broadcast channel $\invecb_3$.
For the values they received from every other player $j \neq i$, they first verify the correctness of the opening $\witness_j$ with respect to $\commitment_j$.
Depending on the output of this check, each participant $i$ does the following:

\begin{itemize}
  \item If the opening $\witness_j$ is valid,
participant $i$ then derives the blinding factor $\blinding_j$ by performing $\blinding_j \gets \enike.\getsharedkey(\witness_j, A_j)$.
  \item If the opening $\witness_j$ is invalid,
the set of participants $\qual$ that followed the protocol honestly cooperate to derive the correct $\blinding_j$ for the cheating participants $\corrupt$.
They do so by broadcasting their respective shares from the misbehaving parties,
and then use the set of received shares to recover $\secretone_j$.

Each remaining participant then performs $\enike.\getsharedkey(\secretone_j, \commitment_j)$,
obtaining $\blinding_j$ as the result.
Due to the assumption that at least $t$ honest participants exist in $\qual$,
this step will always complete successfully.
\end{itemize}

Next, each participant derives $\tweak \gets \avss[\Hash_1].\gettweak(\pubaggregateset, \{\blinding_1, \ldots, \blinding_n \})$,
where  $\pubaggregateset = \{\vsscommit_j \}_{j \in \set{n}}$ is the set of all \AVSS commitments for each participant.
Each participant $i$ then performs $\avss.\aggregatepriv$ and $\avss.\aggregatepub$,
using as input the set of all their received shares $\privaggregateset = \{\share_{ji} \}_{j \in \set{n}}$,
the set of \AVSS commitments $\pubaggregateset$,
and $\tweak$.
Each participant $i$ employs the output of $\avss.\aggregatepriv$ and
$\avss.\aggregatepub$ to determine their aggregated secret key share $\sk_i$ and commitment $\avsscommit$.
The group public key $\PK$ is then derived using $\avsscommit$.
Similarly, the set of each participant's individual public key $\PK_i$ is derived using $\avsscommit$,
and then output as the auxiliary data.

The tweak ensures that the resulting public key is unknown to
an adversary at the time when it picks its contributions to the protocol. As
such, at the beginning of the DKG, the adversary cannot adaptively pick its
contributions so that the resulting public key can be influenced predictably.
In other words, the adversary cannot force the resulting public key to ``hit''
some desired outcome, such as being even or odd, for example.

Participants complete $\finalize$ after setting their status to $\accept$ and
updating their internal state.
The public output at the end of $\dkgkeygen_3$ is $(\pk, \qual, \auxinfo)$.
The private output to each participant $i$ is $(\status, \sk_i)$.



\subsubsection{Recovery Algorithm.}
While most practical DKG applications do not require the participants to actually come together to recover the secret key,
it is possible that some applications may wish to do so. We define the recovery algorithm now.

\begin{itemize}
  \item $\recover(\PK, t, \{ (i, \SK_i) \}_{i \in \coalition}) \rightarrow \SK/\fail$:
    Accepts  the public key $\PK \in \GG$
    and $\lvert \coalition \rvert \geq t$ tuples $ \{ (i, \SK_i) \}_{i \in \coalition}$ consisting of participant $i$'s identifier and their secret key.
    Using these inputs, perform the following steps:
    \begin{enumerate}
      \item Derive $\SK = \avss.\recover(t, \{ (i, \sk_i)\}_{i \in
        \coalition}$).
      \item If $\targetkeygen.\verify(\SK, \PK) \neq 1$, output $\fail$.
      \item Output $\SK$.
    \end{enumerate}
\end{itemize}

\subsection{Security of the Generic Protocol}

We next demonstrate that the generic protocol in Figure~\ref{fg:generic-perfect} is
weakly robust and zero-knowledge.






















%==================================================================================================
